\section{From one primitive plane back to the operator}
\label{sec:operator-ledger}

The exact transform concerns the first mixed trace.  We state the
remaining interfaces before any arithmetic gain can influence the
fixed-shift endpoint.

\subsection{The hierarchy of scopes}

For a possible estimate of
\eqref{eq:primitive-plane-transform}, distinguish:
\begin{enumerate}[label=\textup{(\roman*)},leftmargin=3em]
 \item one dyadic \((U,V)\)-box for one fixed
       \((m,n,\mathfrak a,q)\);
 \item all cofactor boxes for that anchor;
 \item all anchors and erasure keys for one row pair;
 \item all partners incident to one active row;
 \item the complete trace sum; and
 \item a full local-star or operator-norm estimate.
\end{enumerate}
These scopes are not interchangeable.  A small complete trace can result
from cancellation between rows while one local star remains large.  A
fixed-anchor saving can be lost when the number or total mass of anchors
grows.

\subsection{The rectangle is only the linear erasure sector}

The first mixed moment is
\[
 2\tr(ER).
\]
Higher even moments contain one-erasure sectors
\[
 2q\,\tr(ER^{2q-1})
\]
and sectors with two or more erasure edges.  TPC--69 proves that the full
two-erasure sector in every even moment is nonnegative
\citep{WangTPC69}.  Therefore a saving for the rectangle sum does not
bound a complete even moment, and the quadratic sector cannot be used as
a negative correction.

Even a theorem of the form
\[
 \tr(ER)=o\!\left(\sum_{\text{rectangles}}|\cA|\right)
\]
would close only the first signed scalar door.  To reach a native
operator estimate, one must either:
\begin{itemize}
 \item extend primitive-plane control to the local stars and all relevant
       mixed words; or
 \item find a different certificate that converts the same literal
       arithmetic input to \(\|R+E\|\).
\end{itemize}

\subsection{Mass and shorting gates}

Let \(M_X\) be the missing-native map and \(D_M\) its raw atomic row mass.
The normalized missing Gram has the form
\[
 D_M^{\dagger/2}M_X^*M_XD_M^{\dagger/2}
 =I+R+E
\]
on the active row space \citep{WangTPC61,WangTPC68}.  A hypothetical
operator estimate for \(R+E\) still does not finish reassembly.  It must
be combined with:
\begin{enumerate}[label=\textup{(\arabic*)},leftmargin=3em]
 \item a missing-mass comparison in the literal coefficient norm;
 \item a lower-frame estimate for the represented map; and
 \item a negligible nuisance shorting penalty.
\end{enumerate}
In the notation of the existing route, the shorted Gram is
\[
 G_{\rm sh}=G_{\rm dir}-C_N,\qquad C_N\succeq0,
\]
and \(\|C_N\|=o(1)\) is a separate physical statement.  The primitive
cofactor plane neither identifies \(C_N\) nor bounds it.

\subsection{Endpoint exponent ledger}

All identities in this paper cost zero fixed power of \(X\).  Future
arguments may introduce losses for:
\[
 \lambda_{\rm plane},\quad
 \lambda_{\rm anchor},\quad
 \lambda_{\rm native},\quad
 \lambda_{\rm mass},\quad
 \lambda_{\rm rep},\quad
 \lambda_{\rm nuis}.
\]
Only losses proved in a common literal physical scope may be added.  The
total must remain strictly below
\[
 \frac1{400}.
\]
Equality is insufficient for the present endpoint.  A valid theorem
outside this budget may still be valuable, but it cannot be promoted to
the fixed-shift endpoint implication.

\subsection{No fixed-shift arithmetic promotion}

The paper does not establish that the primitive-plane transform is
smaller than its absolute majorant.  Consequently it supplies no new
exponent for the literal physical coefficients, no bound for
\(\|R+E\|\), and no improvement in the endpoint ledger.  Its forward
value is architectural: it closes an exact reindexing interface and
permanently rules out time-ray M\"obius oscillation as the missing
mechanism.
